\documentclass[11pt]{article}
\usepackage{mathblatt}
\makeatletter
\setlength{\mbpfbe}{0mm}% keine Punkte (Bauregel 6.4)
\makeatother
% Rechenraum als Karo (Bauregel 6.9): n Rechenschritte = n mal 10 mm
\newcommand{\karo}[1]{\par\smallskip\noindent\begin{tikzpicture}
  \clip (0,0) rectangle ({\linewidth-1mm},{#1*10mm});
  \draw[black!16,line width=0.3pt,step=5mm] (0,0) grid ({\linewidth},{#1*10mm});
  \end{tikzpicture}\par}
\newcommand{\linie}[1]{\,{\color{black!45}\rule[-1pt]{#1}{0.4pt}}\,}
\newcommand{\pkt}{\textperiodcentered{}}
\begin{document}

\pfheftstil
\pflaufkopf[P10 \pkt{} mehr]{Exponentielles Wachstum \pkt{} Lösungen}{PA50-2}\pfseitevon
\pfheftkopf{Exponentielles Wachstum}{Lösungen \pkt{} mehr}
\begin{pfloesungb}{}
\lzb{1.}{}{2018: 48; 2021: 54,0}{2018: 48; 2021: 54,0}
\end{pfloesungb}
\begin{pfloesungb}{}
\lzb{2.}{}{0: 12 kg; 2: 13,230 kg; 5: $\approx$ 15,315 kg}{0: 12 kg; 2: 13,230 kg; 5: $\approx$ 15,315 kg}
\end{pfloesungb}
\begin{pfloesungb}{}
\lzb{3.}{}{450000 € · 1,06 = 477000 €}{450000 € · 1,06 = 477000 €}
\end{pfloesungb}
\begin{pfloesungb}{}
\lzb{4.}{}{y = 1000 · $0{,}89^{x}$}{y = 1000 · $0{,}89^{x}$}
\end{pfloesungb}
\begin{pfloesungb}{}
\lzb{5.}{}{fallende Strecke durch (0|36), (10|33), (20|30), (30|27), (40|24), (80|12) mit beschrifteter Skala (x-Achse 1 Kästchen = 2 cm)}{fallende Strecke durch (0|36), (10|33), (20|30), (30|27), (40|24), (80|12) mit beschrifteter Skala (x-Achse 1 Kästchen = 5 min, y-Achse 1 Kästchen = 2 cm)}
\end{pfloesungb}
\begin{pfloesungb}{}
\lzb{6.}{}{x: Zeit in Wochen (Anzahl der Wochen)}{12: Anfangsmasse des Ferkels in kg (Woche 0); 1,05: Wachstumsfaktor, Zunahme um 5 \% je Woche; x: Zeit in Wochen (Anzahl der Wochen)}
\end{pfloesungb}
\begin{pfloesungb}{}
\lzb{7.}{}{2 km: $\approx$ 792 hPa; 8 km: $\approx$ 394 hPa}{2 km: $\approx$ 792 hPa; 8 km: $\approx$ 394 hPa}
\end{pfloesungb}
\begin{pfloesungb}{}
\lzb{8.}{}{x-Achse Höhe in km, 2 Kästchen je km; y-Achse Luftdruck in hPa, 1 Kästchen je 50 hPa; sieben Punkte, fallende Kurve}{x-Achse Höhe in km, 2 Kästchen je km; y-Achse Luftdruck in hPa, 1 Kästchen je 50 hPa; sieben Punkte, fallende Kurve}
\end{pfloesungb}
\begin{pfloesungb}{}
\lzb{9.}{}{ja, passt zu 627 hPa: 1000 · 0,89$^{4}$ $\approx$ 627,4 hPa}{ja, passt zu 627 hPa: 1000 · 0,89$^{4}$ $\approx$ 627,4 hPa}
\end{pfloesungb}
\begin{pfloesungb}{}
\lzb{10.}{}{1 h: 5,46 mg; 4,11 mg bei 4 h}{1 h: 5,46 mg; 4,11 mg bei 4 h}
\end{pfloesungb}
\begin{pfloesungb}{}
\lzb{11.}{}{Achsen „Zeit in h“ (1 h je Teilstrich) und „Masse in mg“ (1 mg je Teilstrich), sieben Punkte, glatte fallende Kurve}{Achsen „Zeit in h“ (1 h je Teilstrich) und „Masse in mg“ (1 mg je Teilstrich), sieben Punkte, glatte fallende Kurve}
\end{pfloesungb}
\begin{pfloesungb}{}
\lzb{12.}{}{ $\approx$ 0,62 mg}{$\approx$ 0,62 mg}
\end{pfloesungb}
\begin{pfloesungb}{}
\lzb{13.}{}{Graph 3 passt nicht: beginnt bei 0 kg statt 12 kg}{Graph 1 passt nicht: linear, also jede Woche gleiche Zunahme statt 5 \%; Graph 2 passt: beginnt bei 12 kg, Zuwächse werden größer; Graph 3 passt nicht: beginnt bei 0 kg statt 12 kg}
\end{pfloesungb}
\end{document}